\documentclass{article}

\usepackage{fontspec}
\usepackage{gensymb}
\usepackage{graphicx}
\usepackage{hyperref}
\usepackage{cleveref}
\usepackage{titlesec}
\usepackage{titling}
\usepackage{minted}
\usepackage{natbib}
\usepackage{xcolor}
\usepackage{subcaption}
\setlength{\parindent}{0pt}

\pretitle{\begin{center} \includegraphics[width=200pt]{shitty-logo.png} \\ \vspace{3em} \LARGE}
\title{
  Telekinetic Cowboy Country \\
  \small Organizational structure.
}

\author{
  Emilio Cobos Álvarez \\
  \small \href{mailto:emilio@crisal.io}{emilio@crisal.io}
}

\date{\today}

\begin{document}

\maketitle

\clearpage

\section{Organizational structure}

  Crisal isn't a precisely big corporation, and we tend towards a less
  hierarchical structure. We think excessive hierarchy can have negative impact,
  deteriorating the communication between the customer and the project
  developers, that can be crucial for the success of the project.

  Still, given the scope of {\it Telekinetic Cowboy Country}, we need to have
  some people to direct efforts.

  The idea is that the development effort on our side will be divided as
  follows:

  \begin{itemize}
    \item {\bf Graphics team:} This will be the team leading the rendering
      engine integration and improvement work. This team would be formed by
      our team of experienced Graphics developers, and lead by one of our most
      senior developers. \\

      This team will have little interaction with the final client, and will
      take care of implementing the underlying graphics features that the other
      teams need in order to succeed.

    \item {\bf Game design team:} This team will be responsible of designing the
      different levels, ensure the game is easily playable, and mainly do a lot
      of the playability research and development. \\

      As before, one of the most experienced developers in this area of our
      company will take care of leading this effort.

    \item {\bf Design team:} Well start the game development with a small design
      team lead by one of our designers. This should be enough for the first
      stages of prototyping the game. Over time, we might have the need of
      either contracting more designers, or outsourcing design to other
      third-parties.  \\

      If that's the case, we'll probably need a person taking care of
      relationship with those third-parties, but that's not needed for now.

    \item {\bf Project Manager:} This person will be the responsible of
      coordinating the three aforementioned teams, and coordinating with the
      client.
  \end{itemize}

  Once the prototype is complete, we'll have a {\it Quality Assurance} team that
  will test the new builds on a daily basis and report issues to the teams of
  developers. This team will be formed by people from UBOSOFT. \\

  On the client side of things, we're likely to need to deal with a bunch of
  hierarchy. UBOSOFT has its own organisational structure for the project, but
  we expect the project manager to be able to deal with their team in a
  successful way. If this isn't the case, we may need to form a team just for
  customer relations, but the relationships with the customer so far have been
  quite nice. \\

  There will be bi-weekly progress meetings which will consist of the UBOSOFT
  team, our team leads, our project manager, and (if necessary) other relevant
  developers from our teams. \\

  We expect this organization to be successful, while somewhat unconventional.

\end{document}
